\documentclass[10pt,a4paper]{amsart}
\usepackage{amsmath,amssymb,amsthm,graphicx}
% Retained minimal preamble for the Stacks Project chapter varieties.tex.
% Verbatim excerpt of the chapter preamble: Xy-pic setup (native preamble lines 14-18), the native theorem
% declarations with the mathematical macros (native preamble lines 154-217), and the project own no-op
% definitions of the reference/slogan/history display environments (native wrapper.tex lines 1-3) that the
% counted statement uses for its slogan. Copied mechanically; nothing added, removed or redefined.
% For commutative diagrams we use Xy-pic
\usepackage[all]{xy}

% We use 2cell for 2-commutative diagrams.
\xyoption{2cell}
% Theorem environments.
%
\theoremstyle{plain}
\newtheorem{theorem}[subsection]{Theorem}
\newtheorem{proposition}[subsection]{Proposition}
\newtheorem{lemma}[subsection]{Lemma}

\theoremstyle{definition}
\newtheorem{definition}[subsection]{Definition}
\newtheorem{example}[subsection]{Example}
\newtheorem{exercise}[subsection]{Exercise}
\newtheorem{situation}[subsection]{Situation}

\theoremstyle{remark}
\newtheorem{remark}[subsection]{Remark}
\newtheorem{remarks}[subsection]{Remarks}

\numberwithin{equation}{subsection}

% Macros
%
\def\lim{\mathop{\mathrm{lim}}\nolimits}
\def\colim{\mathop{\mathrm{colim}}\nolimits}
\def\Spec{\mathop{\mathrm{Spec}}}
\def\Hom{\mathop{\mathrm{Hom}}\nolimits}
\def\Ext{\mathop{\mathrm{Ext}}\nolimits}
\def\SheafHom{\mathop{\mathcal{H}\!\mathit{om}}\nolimits}
\def\SheafExt{\mathop{\mathcal{E}\!\mathit{xt}}\nolimits}
\def\Sch{\mathit{Sch}}
\def\Mor{\mathop{\mathrm{Mor}}\nolimits}
\def\Ob{\mathop{\mathrm{Ob}}\nolimits}
\def\Sh{\mathop{\mathit{Sh}}\nolimits}
\def\NL{\mathop{N\!L}\nolimits}
\def\CH{\mathop{\mathrm{CH}}\nolimits}
\def\proetale{{pro\text{-}\acute{e}tale}}
\def\etale{{\acute{e}tale}}
\def\QCoh{\mathit{QCoh}}
\def\Ker{\mathop{\mathrm{Ker}}}
\def\Im{\mathop{\mathrm{Im}}}
\def\Coker{\mathop{\mathrm{Coker}}}
\def\Coim{\mathop{\mathrm{Coim}}}

% Boxtimes
%
\DeclareMathSymbol{\boxtimes}{\mathbin}{AMSa}{"02}

%
% Macros for moduli stacks/spaces
%
\def\QCohstack{\mathcal{QC}\!\mathit{oh}}
\def\Cohstack{\mathcal{C}\!\mathit{oh}}
\def\Spacesstack{\mathcal{S}\!\mathit{paces}}
\def\Quotfunctor{\mathrm{Quot}}
\def\Hilbfunctor{\mathrm{Hilb}}
\def\Curvesstack{\mathcal{C}\!\mathit{urves}}
\def\Polarizedstack{\mathcal{P}\!\mathit{olarized}}
\def\Complexesstack{\mathcal{C}\!\mathit{omplexes}}
% \Pic is the operator that assigns to X its picard group, usage \Pic(X)
% \Picardstack_{X/B} denotes the Picard stack of X over B
% \Picardfunctor_{X/B} denotes the Picard functor of X over B
\def\Pic{\mathop{\mathrm{Pic}}\nolimits}
\def\Picardstack{\mathcal{P}\!\mathit{ic}}
\def\Picardfunctor{\mathrm{Pic}}
\def\Deformationcategory{\mathcal{D}\!\mathit{ef}}
\newenvironment{reference}{}{}
\newenvironment{slogan}{}{}
\newenvironment{history}{}{}
\usepackage{cleveref}
\crefname{originallemma}{original lemma}{original lemmas}
\Crefname{originallemma}{Original lemma}{Original lemmas}
\makeatletter
\newlabel{lemma-power-works}{{0AB3}{source}{}{sourcepointer-1}{}}
\newlabel{lemma-power-works@cref}{{[originallemma][0][]0AB3}{[1][0][]source}{}{}{}}
\makeatother
\crefname{originalreference}{original reference}{original references}
\Crefname{originalreference}{Original reference}{Original references}
\makeatletter
\newlabel{algebra-lemma-integral-closure-Dedekind}{{09IG}{source}{}{sourcepointer-2}{}}
\newlabel{algebra-lemma-integral-closure-Dedekind@cref}{{[originalreference][0][]09IG}{[1][0][]source}{}{}{}}
\makeatother
\crefname{originalreference}{original reference}{original references}
\Crefname{originalreference}{Original reference}{Original references}
\makeatletter
\newlabel{algebra-lemma-integral-under-field}{{00GR}{source}{}{sourcepointer-3}{}}
\newlabel{algebra-lemma-integral-under-field@cref}{{[originalreference][0][]00GR}{[1][0][]source}{}{}{}}
\makeatother
\crefname{originalreference}{original reference}{original references}
\Crefname{originalreference}{Original reference}{Original references}
\makeatletter
\newlabel{algebra-lemma-valuation-ring-cap-field-finite}{{0AAV}{source}{}{sourcepointer-4}{}}
\newlabel{algebra-lemma-valuation-ring-cap-field-finite@cref}{{[originalreference][0][]0AAV}{[1][0][]source}{}{}{}}
\makeatother
\crefname{originallemma}{original lemma}{original lemmas}
\Crefname{originallemma}{Original lemma}{Original lemmas}
\makeatletter
\newlabel{lemma-semi-local-dimension-one-conductor}{{09N2}{source}{}{sourcepointer-5}{}}
\newlabel{lemma-semi-local-dimension-one-conductor@cref}{{[originallemma][0][]09N2}{[1][0][]source}{}{}{}}
\makeatother
\title{Existence of an element of the maximal ideal whose radical is the maximal ideal, matching an element of an integral extension inside the product of localizations}
\author{The Stacks Project Authors}
\date{stacks-project:varieties.tex:local-snapshot-20261011; GNU FDL 1.2 or later}
\begin{document}
\maketitle
\noindent\textit{Editorial note.} All delivered mathematical text is the author's own, copied verbatim from the pinned Stacks Project chapter file varieties.tex; no statement or proof step was rewritten, reordered or removed.\par\medskip
\noindent\textit{Editorial note.} The counted span cites five cross-chapter labels, each recorded in assembly.external\_references with its canonical tag as a navigation pointer.\par\medskip
\noindent\textit{Editorial note.} The author's two parenthetical 'details omitted' remarks concern only the reduction to the reduced case and the closedness of Spec(L) -> Spec(K); the construction itself is written out.\par\medskip
\noindent\textit{Allowed standard prerequisites (editorial).}\begin{itemize}
\item An element with prescribed image in the product of localizations and radical equal to the maximal ideal. Stacks Project, Varieties, Lemma lemma-power-works (tag 0AB7). Allows the reduction to the reduced case at the start of the proof.
\item Integral closure of a one-dimensional Noetherian domain is semi-local Dedekind. Stacks Project, Algebra, Lemma algebra-lemma-integral-closure-Dedekind (tag 034W). Identifies the rings A\_j used in the construction.
\end{itemize}\medskip
\ifdefined\hypertarget\hypertarget{sourcepointer-1}{}\fi
\noindent\textit{Source reference.} Original lemma 0AB3: sources/stacks-chow/tags.txt, source lines 13006--13006. Cross-chapter standard prerequisite cited by the author inside the counted span; it is used as a pointer and its own proof is not part of this item.\par\medskip
\ifdefined\hypertarget\hypertarget{sourcepointer-2}{}\fi
\noindent\textit{Source reference.} Original reference 09IG: sources/stacks-chow/tags.txt, source lines 12039--12039. Cross-chapter standard prerequisite cited by the author inside the counted span; it is used as a pointer and its own proof is not part of this item.\par\medskip
\ifdefined\hypertarget\hypertarget{sourcepointer-3}{}\fi
\noindent\textit{Source reference.} Original reference 00GR: sources/stacks-chow/tags.txt, source lines 617--617. Cross-chapter standard prerequisite cited by the author inside the counted span; it is used as a pointer and its own proof is not part of this item.\par\medskip
\ifdefined\hypertarget\hypertarget{sourcepointer-4}{}\fi
\noindent\textit{Source reference.} Original reference 0AAV: sources/stacks-chow/tags.txt, source lines 12998--12998. Cross-chapter standard prerequisite cited by the author inside the counted span; it is used as a pointer and its own proof is not part of this item.\par\medskip
\ifdefined\hypertarget\hypertarget{sourcepointer-5}{}\fi
\noindent\textit{Source reference.} Original lemma 09N2: sources/stacks-chow/tags.txt, source lines 12200--12200. Cross-chapter standard prerequisite cited by the author inside the counted span; it is used as a pointer and its own proof is not part of this item.\par\medskip
\begin{lemma}
\label{lemma-good-intersection}
Let $A$ be a Noetherian local ring of dimension $1$.
Let $L = \prod A_\mathfrak p$ where the product is over
the minimal primes of $A$. Let $K \to L$ be an integral ring map.
Then there exist $a \in \mathfrak m_A$ and $x \in K$
which map to the same element of $L$ such that $\mathfrak m_A = \sqrt{(a)}$.
\end{lemma}

\begin{proof}
By Lemma \ref{lemma-power-works} we may replace $A$ by
$A/(\bigcap \mathfrak p)$ and assume that $A$ is a reduced ring
(some details omitted).
We may also replace $K$ by the image of $K \to L$.
Then $K$ is a reduced ring. The map $\Spec(L) \to \Spec(K)$ is
surjective and closed (details omitted). Hence $\Spec(K)$ is a finite
discrete space. It follows that $K$ is a finite product of fields.

\medskip\noindent
Let $\mathfrak p_j$, $j = 1, \ldots, m$ be the minimal primes of $A$.
Set $L_j$ be the fraction field of $A_j$ so that
$L = \prod_{j = 1, \ldots, m} L_j$.
Let $A_j$ be the normalization of $A/\mathfrak p_j$. Then
$A_j$ is a semi-local Dedekind domain with at least
one maximal ideal, see
Algebra, Lemma \ref{algebra-lemma-integral-closure-Dedekind}.
Let $n$ be the sum of the numbers of maximal ideals
in $A_1, \ldots, A_m$. For such a maximal ideal $\mathfrak m \subset A_j$
we consider the function
$$
v_{\mathfrak m} : L \to L_j \to \mathbf{Z} \cup \{\infty\}
$$
where the second arrow is the discrete valuation corresponding
to the discrete valuation ring $(A_j)_{\mathfrak m}$ extended
by mapping $0$ to $\infty$. In this way we obtain $n$ functions
$v_1, \ldots, v_n : L \to \mathbf{Z} \cup \{\infty\}$.
We will find an element $x \in K$ such that $v_i(x) < 0$
for all $i = 1, \ldots, n$.

\medskip\noindent
First we claim that for each $i$ there exists an element $x \in K$
with $v_i(x) < 0$. Namely, suppose that $v_i$ corresponds to
$\mathfrak m \subset A_j$. If $v_i(x) \geq 0$ for all $x \in K$,
then $K$ maps into $(A_j)_{\mathfrak m}$ inside the fraction field
$L_j$ of $A_j$.
The image of $K$ in $L_j$ is a field over $L_j$ is
algebraic by Algebra, Lemma \ref{algebra-lemma-integral-under-field}.
Combined we get a contradiction with Algebra, Lemma
\ref{algebra-lemma-valuation-ring-cap-field-finite}.

\medskip\noindent
Suppose we have found an element $x \in K$ such that
$v_1(x) < 0, \ldots, v_r(x) < 0$ for some $r < n$. If $v_{r + 1}(x) < 0$,
then $x$ works for $r + 1$. If not, then choose some $y \in K$ with
$v_{r + 1}(y) < 0$ as is possible by the result of the previous
paragraph. After replacing $x$ by $x^n$ for some $n > 0$,
we may assume $v_i(x) < v_i(y)$ for $i = 1, \ldots, r$. Then
$v_j(x + y) = v_j(x) < 0$ for $j = 1, \ldots, r$ by properties
of valuations and similarly $v_{r + 1}(x + y) = v_{r + 1}(y) < 0$.
Arguing by induction, we find
$x \in K$ with $v_i(x) < 0$ for $i = 1, \ldots, n$.

\medskip\noindent
In particular, the element $x \in K$ has nonzero projection
in each factor of $K$ (recall that $K$ is a finite product of
fields and if some component of $x$ was zero, then one
of the values $v_i(x)$ would be $\infty$). Hence $x$ is
invertible and $x^{-1} \in K$ is an element with
$\infty > v_i(x^{-1}) > 0$ for all $i$. It follows from
Lemma \ref{lemma-semi-local-dimension-one-conductor} that
for some $e < 0$ the element $x^e \in K$ maps to an element of
$\mathfrak m_A/\mathfrak p_j \subset A/\mathfrak p_j$ for all
$j = 1, \ldots, m$. Observe that the cokernel of the map
$\mathfrak m_A \to \prod \mathfrak m_A/\mathfrak p_j$ is
annihilated by a power of $\mathfrak m_A$. Hence after replacing
$e$ by a more negative $e$, we find an element $a \in \mathfrak m_A$
whose image in $\mathfrak m_A/\mathfrak p_j$ is equal to
the image of $x^e$. The pair $(a, x^e)$ satisfies the
conclusions of the lemma.
\end{proof}

\end{document}
